\documentclass[a4paper,twocolumn,landscape]{article}
\usepackage{geometry}
\usepackage{concmath}
\usepackage[OT1]{fontenc}
\usepackage{graphicx}
\usepackage{amsfonts}
\usepackage{amsmath}
\usepackage{calrsfs}
\usepackage{color}
\usepackage[dvipsnames]{xcolor}
\usepackage[many]{tcolorbox}
\usepackage{setspace}              
\usepackage{multicol}  
\geometry{margin=0.5in}
\definecolor{green}{RGB}{52, 235, 168}
\definecolor{violet}{RGB}{165, 83, 237}
\definecolor{lemon}{RGB}{170, 237, 83}
\definecolor{red}{RGB}{240, 125, 84}
\definecolor{blue}{RGB}{105, 255, 242}

\begin{document}
\noindent \colorbox{green}{Exercice 4}\\
1 \begin{align*}
    \vec{v_{R0}} &= \vec{v_G} + \vec{v_{|RG}}\\
    Ec_{|_{R0}} &= \iiint_{(V)} \frac{1}{2}dm.\vec{v(M)}_{|R0}\\
    &= \iiint_{(V)} \frac{1}{2} dm (\vec{v_G} + \vec{v_{|RG}})^2\\
    &= \iiint_{(V)} \frac{1}{2}dm Vc^2 + \iiint_{(V)} \frac{1}{2}dm V_{|RG}^2 + \iiint_{(V)} dm\vec{V_G}.\vec{V_{|RG}}\\
    Ec_{|R0} &= \frac{1}{2}mV_G^2 + Ec_{|RG} + \vec{V_G} \iiint_{(V)} dm\vec{V_{|RG}}\\
    &\vec{v}(M)_{|RG} = \frac{d}{dt} \vec{GM}|_{RG}\\
    &\vec{V_G} \iiint_{(V)} dm\vec{V_{|RG}} = \vec{v_G}\frac{d}{dt}(\iiint_{(V)} \vec{GM}|_{RG}.dm) = 0
\end{align*}
Donc \fbox{\(Ec_{R0} = \frac{1}{2}mV_{G}^2 + Ec|_{R0}\)}\\
2 \begin{align*}
    J = mR^2\\
\end{align*}
1a vitesse \(\omega\)\begin{align*}
    Ec &= \frac{1}{2}J\omega^2\\
    &= mR^2\omega^2 
\end{align*}

3 Donc \begin{align*}
    Ec &= \frac{1}{2}J\omega^2 + \frac{1}{2}mV_G^2\\
    &= \frac{1}{2}mR^2\omega^2 + \frac{1}{2}mR^2\omega^2\\
    Ec_{|_{R0}} &= mr^2\omega^2
\end{align*}
4 Par conservation de Em : 
\(mV_0^2 = mgsin\alpha\)
d = \(\frac{v_0^2}{gsin\alpha}\)

\noindent \colorbox{green}{Exercice 5}\\
\begin{equation*}
    \vec{Fie} = -m\vec{a_e} = -m\vec{a}
\end{equation*}
\underline{À l'équilibre} dans R0 : \(\vec{0} = \vec{P} + \vec{T} + \vec{Fie}\)\\
\begin{align*}
    tan\theta_e &= \frac{||\vec{Fie}||}{||\vec{P}||}\\
    tan\theta_e &= \frac{a}{g}
\end{align*}
2. PFD à A dans R (non gal)\begin{align*}
    m\frac{d\vec{v}}{dt}|_{R0} &= \vec{P} + \vec{Fie} + \vec{T} &\vec{v_A}|_{R} = l\dot{\theta}\vec{u_\theta}\\
    &ml\ddot{\theta}\vec{u_\theta} - ml\dot{\theta}^2\vec{u_r} = -mg \vec{u_z} - ma\vec{u_y} - T\vec{u_r} &\frac{d\vec{v_A}}{dt}|_{R} = l\ddot{\theta}\vec{u_\theta} - l\dot{\theta}^2\vec{u_r}\\
    &ml\ddot{\theta} + 0 = -mg sin\theta + macos\theta + 0 \text{ projeter sur \(\vec{u_\theta}\) }\\
    &l\ddot{\theta} + 0 = -g sin\theta + acos\theta \\
    & \ddot{\theta} + \frac{g}{l}sin\theta - \frac{a}{l}cos\theta = 0 \indent(1)& \theta = \theta_e + \varepsilon, \varepsilon << 1
\end{align*}
3. d.l. ordre 1 en \(\varepsilon\) \begin{align*}
    sin(\theta_e + \varepsilon) = sin\theta_e + \varepsilon cos\theta_e\\
    cos(\theta_e + \varepsilon) = cos\theta_e - \varepsilon sin\theta_e
\end{align*}
\begin{align*}
    &(1) \indent \ddot{\varepsilon} + \frac{g}{l}sin\theta_e + \frac{g}{l}cos\theta_e - \frac{a}{l}cos\theta_e + \frac{a}{l}sin\theta_e.\varepsilon = 0\\
    &\ddot{\varepsilon} + (\frac{g}{l}cos\theta_e + \frac{a}{l}sin\theta_e).\varepsilon(t) = 0
\end{align*}
oscillateur harmonique \(\Rightarrow\)
\(\varepsilon\) varie sinuisoïdalement autour de 0 et \(\theta\) varie sinuisoïdalement autour de \(\theta_e\)\\
Implique que l'équilibre est stable\\
*\(T_0 = \frac{2\pi}{\omega_0}\)
\(\omega_0^2 = \frac{g}{l}cos\theta_e + \frac{a}{l}sin\theta_e\)
\(\frac{1}{1+tan^2\theta_e} = cos^2\theta_e = \frac{1}{1+\frac{a^2}{g^2}} = \frac{g^2}{a^2 + g^2}\)\\
\(sin^2\theta_e = \frac{a^2}{a^2+g^2}\)
\begin{eqnarray*}
    \omega_0^2 = \frac{g^2}{l}\frac{1}{\sqrt{a^2+g^2}} + \frac{a^2}{l}\frac{1}{\sqrt{a^2 + g^2}}\\
    \omega_0^2 = \frac{1}{l}\sqrt{a^2+g^2}
\end{eqnarray*}

\noindent \colorbox{green}{Exercice 8}\\
Dans R : réf lié à la tige en rotation autour de Oz dans R0 galiléen\begin{align*}
    \vec{Fie} &= m\omega^2 \vec{HM}\\
    \vec{Fie} &= m\omega^2 x sin\alpha \vec{u_r}
\end{align*}
\underline{À l'équilibre} \(\vec{0} = \vec{P} + \vec{R} + \vec{Fie}\)\\
\begin{align*}
    tan \alpha &= \frac{||\vec{P}||}{||\vec{Fie}||}\\
    &= \frac{mg}{m\omega^2x_esin\alpha}\\
    \Rightarrow x_e &= \frac{gcos\alpha}{\omega^2 sin\alpha} \indent (OK)
\end{align*}
On va dire \(x = x_e + \varepsilon\)\\
PFD dans R à l'anneau\begin{align*}
    m\vec{a}(M)|_R &= \vec{P} + \vec{R} + \vec{Fie} + \vec{Fic}\\
    m\ddot{x}\vec{u_x} &= \vec{P}+\vec{R}+\vec{Fie} -2m\vec{v_r}\wedge \vec{\omega}\\
    \text{ Tout projeter selon \(\vec{u_x}\)}
\end{align*}
\end{document}